% Deliberately broken deck: every frame violates an econ-slides rule.
% Negative and review fixtures for check_deck.py, plus explicit non-violations.
\documentclass[11pt, aspectratio=169]{beamer}
\usepackage{econ-slides-house}

\begin{document}

% VIOLATION 1: frame title wraps to two lines
\begin{frame}{This frame title is far too long to fit on a single line and therefore wraps onto a second line which the skill forbids}
  \begin{itemize}
    \item Short bullet
  \end{itemize}
\end{frame}

% VIOLATION 2: wrapped bullet + too many bullets
\begin{frame}{Too many bullets and one that wraps}
  \begin{itemize}
    \item This bullet goes on and on with clause after clause after clause so that it cannot possibly fit on one line at this font size and must wrap
    \item Second bullet
    \item Third bullet
    \item Fourth bullet
    \item Fifth bullet
    \item Sixth bullet
    \item Seventh bullet
  \end{itemize}
\end{frame}

% VIOLATION 3: text pushed past the right page edge
\begin{frame}{Edge overflow}
  \noindent\hspace{0.8\textwidth}\mbox{this text is pushed far beyond the right edge of the slide}
\end{frame}

% VIOLATION 4: overlay as decoration — a \pause chain
\begin{frame}{Pause chain}
  \begin{itemize}
    \item One \pause
    \item Two \pause
    \item Three \pause
    \item Four
  \end{itemize}
\end{frame}

% VIOLATION 5: wrapped enumerate item (the punchline-slide failure mode)
\begin{frame}{Wrapped numbered contribution}
  \begin{enumerate}
    \item This numbered contribution rambles through clause after clause after clause until it cannot fit on one line and wraps
    \item Short item
  \end{enumerate}
\end{frame}

% NON-VIOLATION: an em-dash filling an empty table cell must NOT be flagged
% as a bullet (false-positive regression test)
\begin{frame}{Em-dash table cell is not a bullet}
  \centering
  \begin{tabular}{lcc}
    \toprule
    & (1) & (2) \\
    \midrule
    Effect & 0.12 & --- \\
    Control mean & --- & 0.48 \\
    \bottomrule
  \end{tabular}
\end{frame}

% VIOLATION 6: navigation target does not exist
\begin{frame}{Broken navigation}
  \hyperlink{app:missing}{\beamergotobutton{Missing backup}}
\end{frame}

% VIOLATION 7: two generic punchline treatments compete on one frame
\begin{frame}{Too many highlights}
  \KeyIdea{First key idea} and \KeyIdea{second key idea} cannot both dominate.
\end{frame}

% VIOLATION 8: a result box and a KeyIdea compete as primary emphasis
\begin{frame}{A box and a key phrase cannot both dominate}
  \KeyIdea{First primary treatment}
  \begin{ResultBox}[Second primary treatment]
    One frame still needs one visual hierarchy.
  \end{ResultBox}
\end{frame}

% REVIEW 1: prose-only columns should prompt a single-canvas review
\begin{frame}{Text columns need a real comparison}
  \begin{columns}
    \begin{column}{0.48\textwidth}Ordinary prose on the left.\end{column}
    \begin{column}{0.48\textwidth}Ordinary prose on the right.\end{column}
  \end{columns}
\end{frame}

% REVIEW 2: paired minipages are also a prose-only two-column layout
\begin{frame}{Text minipages need a real comparison}
  \begin{minipage}{0.48\textwidth}Ordinary prose on the left.\end{minipage}\hfill
  \begin{minipage}{0.48\textwidth}Ordinary prose on the right.\end{minipage}
\end{frame}

% VIOLATION 9: arrow-shaped list marker
\begin{frame}{Arrow bullets flatten the hierarchy}
  \begin{itemize}
    \item[$\Rightarrow$] This should be an ordinary bullet or interpretation line
    \item[→] Unicode arrows are equally mechanical
  \end{itemize}
\end{frame}

% VIOLATION 10: prohibited opening label
\begin{frame}{The economic question}
  This structural object must be called the Key question.
\end{frame}

% REVIEW: a concise question belongs in the frame title, not its subtitle
\begin{frame}{Key question}
  \framesubtitle{Can a \textcolor{cAccentA}{short title} state the question directly?}
  The subtitle hides the opening object.
\end{frame}

% NON-VIOLATION: setup may accompany a question written in the body
\begin{frame}{Key question}
  Can a short title state the question directly?
\end{frame}

% REVIEW: a roadmap is an orientation device, not an emphasis hierarchy
\begin{frame}{Roadmap}
  \begin{enumerate}
    \item \textbf{Setting and data}
    \item \textbf{Design and result}
    \item \textbf{Interpretation}
  \end{enumerate}
\end{frame}

% NON-VIOLATION: regular-weight Roadmap entries are the default
\begin{frame}{Roadmap}
  \begin{enumerate}
    \item Setting and data
    \item Design and result
    \item Interpretation
  \end{enumerate}
\end{frame}

% STRUCTURE CONTROL: the checker does not impose a Key-question/Roadmap order.
\begin{frame}
  \frametitle{This Paper}
  The answer appears early.
\end{frame}
\begin{frame}{Data appears too soon}
  The paper and audience determine whether a Roadmap is useful here.
\end{frame}

% REVIEW 3: a static exhibit needs a bottom interpretation line
\begin{frame}{A table without an economic reading}
  \centering
  \begin{tabular}{lc}
    \toprule
    Treatment & 0.12 \\
    \bottomrule
  \end{tabular}
\end{frame}

% VIOLATION 12: the conclusion is not a navigation hub
\begin{frame}{Conclusion}
  \begin{itemize}
    \item Main answer
    \item First support
    \item Second support
    \item Robustness inventory
    \item New detail
  \end{itemize}
  \PlaceNav{\hyperlink{app:existing}{\beamergotobutton{Backup}}}
\end{frame}
\begin{frame}{Defined target}
  \hypertarget{app:existing}{}Target exists so this is not a dead-link test.
\end{frame}

\end{document}
